\chapter{盡心章句上}
\kaishi*

\section{盡心知性}

孟子曰：「盡其心者，知其性也。知其性，則知天矣。」\jiazhu[1]{性有仁義禮智之端，心以制之，惟心爲正。人能盡極其心以思行善，則可謂知其性矣。知其性，則知天道之貴善者也。}

「存其心，養其性，所以事天也。」\jiazhu[1]{能存其心，養育其正性，可謂仁人。天道好生，仁人亦好生；天道無親，惟仁是與。行與天合，故曰所以事天。}

「夭壽不貳，修身以俟之，所以立命也。」\jiazhu[1]{貳，二也。仁人之行，一度而已，雖見前人或夭或壽，終無二心改易其道。夭若顏淵，壽若邵公，皆歸之命。修正其身以待天命，此所以立命之本也。}\par \jiazhu[1]{章指言：盡心竭性，足以承天。夭壽禍福，秉心不違。立命之道，惟是爲珍。}

\section{知命者不立巖牆之下}

孟子曰：「莫非命也，順受其正。」\jiazhu[1]{莫，無也。人之終無非命也。命有三名：行善得善曰受命，行善得惡曰遭命，行惡得惡曰隨命。惟順受命，爲受其正也。}

「是故知命者，不立乎巖牆之下。盡其道而死者，正命也。」\jiazhu[1]{知命者欲趨於正，故不立巖牆之下，恐壓覆也。盡修身之道，以壽終者，爲得正命也。}

「桎梏死者，非正命也。」\jiazhu[1]{畏壓溺，禮所不弔，故曰非正命也。}\par \jiazhu[1]{章指言：人必趨命，貴受其正。巖牆之疑，君子遠之。}

\section{求則得之，舍則失之}

孟子曰：「求則得之，舍則失之，是求有益於得也，求在我者也。」\jiazhu[1]{謂脩仁行義，事在於我。我求則得，我舍則失，故求有益於得也。}

「求之有道，得之有命，是求無益於得也，求在外者也。」\jiazhu[1]{謂賢者脩其天爵而人爵從之，故曰求之有道也。脩天爵者或得或否，故言得之有命也。爵禄須知己，知己者在外，非身所專，是以云求無益於得也，求在外者也。}\par \jiazhu[1]{章指言：爲仁由己，富貴在天。故孔子曰：「如不可求，從吾所好。」}

\section{萬物皆備於我}

孟子曰：「萬物皆備於我矣。反身而誠，樂莫大焉。」\jiazhu[1]{物，事也；我，身也。普謂人爲成人已往，皆備知天下萬物常有所行矣。誠者，實也。反自思其身所施行，能皆實而無虛，則樂莫大焉。}

「強恕而行，求仁莫近焉。」\jiazhu[1]{當自強勉以忠恕之道，求仁之術，此最爲近。}\par \jiazhu[1]{章指言：每必以誠，恕已而行，樂在其中，仁之至也。}

\section{行之不著，習矣不察}

孟子曰：「行之而不著焉，習矣而不察焉，終身由之而不知其道者，衆也。」\jiazhu[1]{人皆有仁義之心，日自行之於其所愛，而不能著明其道以施於大事。仁妻愛子，亦以習矣，而不能察知可推以爲善也。由，用也。終身用之，以爲自然，不究其道可成君子。此衆庶之人也。}\par \jiazhu[1]{章指言：人有仁端，達之以爲道。凡夫用之，不知其爲寶也。}

\section{人不可以無恥}

孟子曰：「人不可以無恥。」\jiazhu[1]{人不可以無所羞恥也。《論語》曰：「行己有恥。」}

「無恥之恥，無恥矣。」\jiazhu[1]{人能恥己之無所恥，是爲改行從善之人，終身無復有恥辱之累也。}\par \jiazhu[1]{章指言：恥身無分，獨無所恥，斯必遠辱，不爲憂矣。}

\section{恥之於人大矣}

孟子曰：「恥之於人大矣。爲機變之巧者，無所用恥焉。」\jiazhu[1]{恥者，爲不正之道，正人之所恥爲也。今造機變阱陷之巧以攻戰者，非古之正道也，取爲一切可勝敵也，宜無以錯於廉恥之心也。}

「不恥不若人，何若人有？」\jiazhu[1]{不恥不如古之聖人，何有如賢人之名也？}\par \jiazhu[1]{章指言：不慕大人，何能有恥？是以隰朋愧不及黃帝，佐齊桓以有勳；顏淵慕虞舜，仲尼歎庶幾之云。}

\section{古之賢王好善忘勢}

孟子曰：「古之賢王好善而忘勢。」\jiazhu[1]{樂善而自卑，若高宗得傅說而稟命。}

「古之賢士何獨不然？樂其道而忘人之勢。」\jiazhu[1]{何獨不然？何獨不有所樂有所忘也？樂道守志，若許由洗耳，可謂忘人之勢矣。}

「故王公不致敬盡禮，則不得亟見之。見且由不得亟，而況得而臣之乎？」\jiazhu[1]{亟，數也。若伯夷非其君不事，伊尹樂堯舜之道，不致敬盡禮，可數見之乎？作者七人，隱各有方，豈可得而臣之？}\par \jiazhu[1]{章指言：王公尊賢，以貴下賤之義也；樂道忘勢，不以富貴動心之分也。各崇所尚，則義不虧矣。}

\section{孟子謂宋句踐}

孟子謂宋句踐曰：「子好遊乎？吾語子遊。人知之，亦囂囂；人不知，亦囂囂。」\jiazhu[1]{宋，姓也；句踐，名也。好以道德遊，欲行其道者。囂囂，自得無欲之貌。}

曰：「何如斯可以囂囂矣？」\jiazhu[1]{句踐問何執守可囂囂也。}

曰：「尊德樂義，則可以囂囂矣。」\jiazhu[1]{尊，貴也。孟子曰：能貴德而履之，樂義而行之，則可以囂囂無欲矣。}

「故士窮不失義，達不離道。窮不失義，故士得己焉；達不離道，故民不失望焉。」\jiazhu[1]{窮不失義，不爲不義而苟得，故得己之本性也。達不離道，思利民之道，故民不失其望也。}

「古之人，得志，澤加於民；不得志，脩身見於世。窮則獨善其身，達則兼善天下。」\jiazhu[1]{古之人得志君國，則德澤加於民人；不得志，謂賢者不遭遇也。見，立也。獨治其身以立於世間，不失其操也。是故獨善其身。達，謂得行其道，故能兼善天下也。}\par \jiazhu[1]{章指言：內定常滿，囂囂無憂，可出可處，故云以遊。修身立世，賤不失道，達善天下，乃用其寶。句踐好遊，未得其要，孟子言之，然後乃喻。}

\section{待文王而後興者凡民也}

孟子曰：「待文王而後興者，凡民也。若夫豪傑之士，雖無文王，猶興。」\jiazhu[1]{凡民，無異知者也，故須文王之大化，乃能自興起以趨善道。若夫豪傑，才知千萬於凡人者，雖不遭文王，猶能自起以善，守身正行，不陷溺也。}\par \jiazhu[1]{章指言：小人待化，乃不辟邪；君子特立，不爲俗移，故稱豪傑自興也。}

\section{附之以韓魏之家}

孟子曰：「附之以韓魏之家，如其自視欿然，則過人遠矣。」\jiazhu[1]{附，益也。韓、魏，晉六卿之富者也。言人既自有家，復益以韓魏百乘之家，其富貴已美矣。而其人欿然不足，自知仁義之道不足也，此則過人甚遠矣。}\par \jiazhu[1]{章指言：人情富盛，莫不驕矜，若能欿然謂不如人，非但免過，卓絕乎凡也。}

\section{佚道使民，生道殺民}

孟子曰：「以佚道使民，雖勞不怨。」\jiazhu[1]{謂教民趨農，役有常時，不使失業。當時雖勞，後獲其利則佚矣。若亟其乘屋之類也，故曰不怨。}

「以生道殺民，雖死不怨殺者。」\jiazhu[1]{謂殺大辟之罪者，以坐殺人故也。殺此罪人者，其意欲生民也，故雖伏罪而死，不怨殺者。}\par \jiazhu[1]{章指言：勞人欲以佚之，殺人欲以生之，則民無怨讟也。}

\section{霸者之民與王者之民}

孟子曰：「霸者之民，驩虞如也；王者之民，皞皞如也。殺之而不怨，利之而不庸，民日遷善而不知爲之者。」\jiazhu[1]{霸者行善恤民，恩澤暴見易知，故民驩虞樂之也。王者道大法天，浩浩而德難見也。殺之不怨，故曰殺人而不怨也。庸，功也。利之使趨時而農，六畜繁息，無凍餓之老，而民不知獨是王者之功。修其庠序之教，使日遷善，亦不能覺知誰爲之者，言化大也。}

「夫君子所過者化，所存者神，上下與天地同流，豈曰小補之哉？」\jiazhu[1]{君子通於聖人，聖人如天。過此世能化之，存在此國，其化如神。故言與天地同流也。天地化物，歲成其功，豈曰使成人知其小補益也？}\par \jiazhu[1]{章指言：王政浩浩，與天地同道；霸者德小，民人速睹。是以賢者志其大者也。}

\section{仁言不如仁聲}

孟子曰：「仁言不如仁聲之入人深也。」\jiazhu[1]{仁言，政教法度之言也。仁聲，樂聲雅頌也。仁言之政雖明，不如雅頌感人心之深也。}

「善政不如善教之得民也。善政，民畏之；善教，民愛之。善政得民財，善教得民心。」\jiazhu[1]{善政使民不違上，善教使民尚仁義，心易得也。畏之不逋怠，故賦役舉而財聚於一家也。愛之樂風化而上下親，故歡心可得也。}\par \jiazhu[1]{章指言：明法審令，民趨君命；崇寬務化，民愛君德。故曰移風易俗，莫善於樂。}

\section{不學而能，不慮而知}

孟子曰：「人之所不學而能者，其良能也；所不慮而知者，其良知也。」\jiazhu[1]{不學而能，性所自能。良，甚也，是人之所能甚也。知亦猶是能也。}

「孩提之童，無不知愛其親者；及其長也，無不知敬其兄也。」\jiazhu[1]{孩提，二三歲之間在襁褓，知孩笑可提抱者也。少知愛親，長知敬兄，此所謂良能良知也。}

「親親，仁也；敬長，義也。無他，達之天下也。」\jiazhu[1]{人仁義之心少而皆有之，欲爲善者無他，達，通也，但通此親親敬長之心，施之天下人而已。}\par \jiazhu[1]{章指言：本性良能，仁義是也。達之天下，恕乎己也。}

\section{舜居深山}

孟子曰：「舜之居深山之中，與木石居，與鹿豕遊，其所以異於深山之野人者幾希。」\jiazhu[1]{舜耕歷山之時，居木石之間，鹿豕近人，若與人遊也。希，遠也。當此之時，舜與野人相去豈遠哉？}

「及其聞一善言，見一善行，若決江河，沛然莫之能禦也。」\jiazhu[1]{舜雖外與野人同其居處，聞人一善言則從之，見人一善行則識之，沛然不疑，若江河之流，無能禦止其所欲行。}\par \jiazhu[1]{章指言：聖人潛隱，辟若神龍，亦能飛天，亦能小同，舜之謂也。}

\section{無爲其所不爲}

孟子曰：「無爲其所不爲，無欲其所不欲，如此而已矣。」\jiazhu[1]{無使人爲己所不欲爲者，無使人欲己之所不欲者，每以身況之，如此則人道足也。}\par \jiazhu[1]{章指言：己所不欲，勿施於人，仲尼之道也。}

\section{德慧術知存乎疢疾}

孟子曰：「人之有德慧術知者，恒存乎\zhuyin{疢疾}。」\jiazhu[1]{人所以有德行智慧道術才智者，在於有疢疾之人。疢疾之人又力學，故能成德。}

「獨孤臣孽子，其操心也危，其慮患也深，故達。」\jiazhu[1]{此即人之疢疾也。自以孤微，懼於危殆之患，而深慮之，勉爲仁義，故至於達也。}\par \jiazhu[1]{章指言：孤孽自危，故能顯達；膏粱難正，多用沈溺。是故在上不驕，以戒諸侯也。}

\section{事君者四科}

孟子曰：「有事君人者，事是君則爲容悅者也。」\jiazhu[1]{事君求君之意，爲苟容以悅君而已。}

「有安社稷臣者，以安社稷爲悅者也。」\jiazhu[1]{忠臣志在安社稷而後悅也。}

「有天民者，達可行於天下而後行之者也。」\jiazhu[1]{天民，知道者也，可行而行，可止而止。}

「有大人者，正己而物正者也。」\jiazhu[1]{大人，大丈夫不爲利害動移者也。正己物正，象天不言而萬物化成也。}\par \jiazhu[1]{章指言：容悅凡臣，社稷股肱，天民行道，大人正身。凡此四科，優劣之差。}

\section{君子有三樂}

孟子曰：「君子有三樂，而王天下不與存焉。父母俱存，兄弟無故，一樂也；仰不愧於天，俯不怍於人，二樂也；得天下英才而教育之，三樂也。」\jiazhu[1]{天下之樂不得與此三樂之中。兄弟無故，無他故。不愧天，又不怍人，心正無邪也。育，養也。教養英才，成之以道，皆樂也。}

「君子有三樂，而王天下不與存焉。」\jiazhu[1]{孟子重言，是美之也。}\par \jiazhu[1]{章指言：保親之養，兄弟無他，誠不愧天，育養英才。賢人能之，樂過萬乘，孟子重焉，一章再云也。}

\section{廣土衆民非君子所性}

孟子曰：「廣土衆民，君子欲之，所樂不存焉。」\jiazhu[1]{廣土衆民，大國諸侯也。所樂不存，樂行禮也。}

「中天下而立，定四海之民，君子樂之，所性不存焉。」\jiazhu[1]{中天下而立，謂王者。所性不存，謂性仁義也。}

「君子所性，雖大行不加焉，雖窮居不損焉，分定故也。」\jiazhu[1]{大行，行政於天下。窮居不失性也，分定故不變。}

「君子所性，仁義禮智根於心。其生色也，睟然見於面，盎於背，施於四體。四體不言而喻。」\jiazhu[1]{四者根生於心，色見於面。睟然，潤澤之貌也。盎，視其背而可知，其背盎盎然，盛流於四體。四體有匡國之綱，雖口不言，人以曉喻而知之也。}\par \jiazhu[1]{章指言：臨蒞天下，君國子民，君子之樂，尚不與存。仁義內充，身體履方，四支不言，蟠辟用張。心邪意溺，進退無容，於是之際，知其不同也。}

\section{文王善養老}

孟子曰：「伯夷辟紂，居北海之濱，聞文王作興，曰：『盍歸乎來！吾聞西伯善養老者。』太公辟紂，居東海之濱，聞文王作興，曰：『盍歸乎來！吾聞西伯善養老者。』」\jiazhu[1]{已說於上篇。}

「天下有善養老，則仁人以爲己歸矣。」\jiazhu[1]{天下有能若文王者，仁人將復歸之矣。}

「五畝之宅，樹牆下以桑，匹婦蠶之，則老者足以衣帛矣。五母雞，二母彘，無失其時，老者足以無失肉矣。百畝之田，匹夫耕之，八口之家足以無飢矣。」\jiazhu[1]{五雞二彘，八口之家畜之，足以爲畜產之本也。}

「所謂西伯善養老者，制其田里，教之樹畜，導其妻子，使養其老。五十非帛不煖，七十非肉不飽。不煖不飽，謂之凍餒。文王之民無凍餒之老者，此之謂也。」\jiazhu[1]{所謂無凍餒者，教導之使可以養老者耳，非家賜而人益之也。}\par \jiazhu[1]{章指言：王政普大，教其常業，各養其老，使不凍餒。二老聞之，歸身自託。衆鳥不羅，翔鳳來集，亦斯類也。}

\section{易其田疇，民可使富}

孟子曰：「易其田疇，薄其稅斂，民可使富也。」\jiazhu[1]{易，治也。疇，一井也。教民治其田疇，薄其稅斂，不踰什一，則民富矣。}

「食之以時，用之以禮，財不可勝用也。」\jiazhu[1]{食取其征賦以時，用之以常禮，不踰禮以費財也。故畜積有餘，財不可勝用也。}

「民非水火不生活。昏暮叩人之門户求水火，無弗與者，至足矣。聖人治天下，使有菽粟如水火。菽粟如水火，而民焉有不仁者乎？」\jiazhu[1]{水火能生人，有不愛者，至饒足故也。菽粟饒多若是，民皆輕施於人，何有不仁者也？}\par \jiazhu[1]{章指言：教民之道，富而節用，畜積有餘，焉有不仁？故曰倉廩實知禮節也。}

\section{登東山而小魯}

孟子曰：「孔子登東山而小魯，登太山而小天下。故觀於海者難爲水，遊於聖人之門者難爲言。」\jiazhu[1]{所覽大者意大，觀小者志小也。}

「觀水有術，必觀其瀾。」\jiazhu[1]{瀾，水中大波也。}

「日月有明，容光必照焉。」\jiazhu[1]{容光，小郤也。言大明照幽微也。}

「流水之爲物也，不盈科不行；君子之志於道也，不成章不達。」\jiazhu[1]{盈，滿也；科，坎也。流水滿坎乃行，以喻君子學必成章乃仕進也。}\par \jiazhu[1]{章指言：弘大明者無不照，包聖道者成其仁。是故賢者志大，宜爲君子。}

\section{雞鳴而起}

孟子曰：「雞鳴而起，孳孳爲善者，舜之徒也。雞鳴而起，孳孳爲利者，蹠之徒也。欲知舜與蹠之分，無他，利與善之間也。」\jiazhu[1]{蹠，盜蹠也。舜、蹠之分，以此別之。}\par \jiazhu[1]{章指言：好善從舜，好利從蹠，明明求之，常若不足。君子小人，各一趣也。}

\section{楊墨之道}

孟子曰：「楊子取爲我，拔一毛而利天下，不爲也。」\jiazhu[1]{楊子，楊朱也。爲我，爲己也。拔己一毛以利天下之民，不爲也。}

「墨子兼愛，摩頂放踵利天下，爲之。」\jiazhu[1]{墨子，墨翟也。兼愛他人。摩突其頂，下至於踵，以利天下，己樂爲之也。}

「子莫執中。」\jiazhu[1]{子莫，魯之賢人也，其性中和專一者也。}

「執中爲近之。執中無權，猶執一也。」\jiazhu[1]{執中和近聖人之道，然不權。聖人之重權，執中而不知權，猶執一介之人，不得時變也。}

「所惡執一者，爲其賊道也，舉一而廢百也。」\jiazhu[1]{所以惡執一者，爲其不知權，以一知而廢百道也。}\par \jiazhu[1]{章指言：楊墨放蕩，子莫執一，聖人量時，不取此術。孔子行止，唯義所在。}

\section{飢渴害心}

孟子曰：「飢者甘食，渴者甘飲，是未得飲食之正也，飢渴害之也。」\jiazhu[1]{飢渴害其本所以知味之性，令人強甘之。}

「豈惟口腹有飢渴之害？人心亦皆有害。」\jiazhu[1]{爲利欲所害，亦猶飢渴得之。}

「人能無以飢渴之害爲心害，則不及人不爲憂矣。」\jiazhu[1]{人能守正，不爲邪利所害，雖謂富貴之事不及逮人，猶爲君子，不爲善人所憂患也。}\par \jiazhu[1]{章指言：飢不妄食，忍情抑欲，賤不失道，不爲苟求。能無心害，夫將何憂？}

\section{柳下惠之介}

孟子曰：「柳下惠不以三公易其介。」\jiazhu[1]{介，大也。柳下惠執弘大之志，不恥汙君，不以三公榮位易其大量也。}\par \jiazhu[1]{章指言：柳下惠不恭，用志大也。無可無否，以賤爲貴也。}

\section{有爲者譬若掘井}

孟子曰：「有爲者辟若掘井，掘井九軔而不及泉，猶爲棄井也。」\jiazhu[1]{有爲爲仁義也。軔，八尺也。雖深而不及泉，喻有爲者中道而盡棄前行也。}\par \jiazhu[1]{章指言：爲仁由己，必在究之，九軔而輟，無益成功。論之一簣，義與此同。}

\section{堯舜性之，湯武身之}

孟子曰：「堯舜，性之也；湯武，身之也；五霸，假之也。」\jiazhu[1]{性之，性好仁，自然也。身之，體之行仁，視之若身也。假之，假仁以正諸侯也。}

「久假而不歸，惡知其非有也？」\jiazhu[1]{五霸若能久假仁義，譬若假物久而不歸，安知其不真有也？}\par \jiazhu[1]{章指言：仁在性體，其次假借，用而不已，實何以易？在其勉之也。}

\section{伊尹放太甲}

公孫丑曰：「伊尹曰：『予不狎于不順。』放太甲于桐，民大悅。太甲賢，又反之，民大悅。賢者之爲人臣也，其君不賢，則固可放與？」\jiazhu[1]{丑怪伊尹賢者而放其君，何也？}

孟子曰：「有伊尹之志，則可；無伊尹之志，則篡也。」\jiazhu[1]{大臣秉忠，志若伊尹欲寧殷國，則可放惡而不即立君，宿留冀改而復之。如無伊尹之忠，見間乘利，篡心乃生，何可放也？}\par \jiazhu[1]{章指言：憂國忘家，意在出身，志在寧君，放惡攝政，伊周有焉。凡人志異，則生篡心也。}

\section{君子不耕而食}

公孫丑曰：「《詩》曰：『不素餐兮。』君子之不耕而食，何也？」\jiazhu[1]{《詩·魏風·伐檀》之篇也。無功而食謂之素餐。世之君子有不耕而食者，何也？}

孟子曰：「君子居是國也，其君用之，則安富尊榮；其子弟從之，則孝悌忠信。『不素餐兮』，孰大於是？」\jiazhu[1]{君子能使人化其道德，移其習俗，君安國富而保其尊榮，子弟孝悌而樂忠信。不素餐之功，誰大於是？何爲不可以食祿？}\par \jiazhu[1]{章指言：君子正己以立於世，世美其道，君臣是貴，所過者化，何素餐之謂也？}

\section{王子墊問士何事}

王子墊問曰：「士何事？」\jiazhu[1]{齊王子名墊也。問士當何事爲事也？}

孟子曰：「尚志。」\jiazhu[1]{尚，上也。士當貴上於用志也。}

曰：「何謂尚志？」曰：「仁義而已矣。殺一無罪，非仁也；非其有而取之，非義也。居惡在？仁是也。路惡在？義是也。居仁由義，大人之事備矣。」\jiazhu[1]{孟子言志之所尚，仁義而已矣。不殺無罪，不取非有者，爲仁義。欲知其所當居者，仁爲上；所由者，義爲貴。大人之事備也。}\par \jiazhu[1]{章指言：人當尚志，志於善也。善之所由，仁與義也。欲使王子無過差也。}

\section{仲子之義}

孟子曰：「仲子，不義與之齊國而弗受，人皆信之，是舍簞食豆羹之義也。」\jiazhu[1]{仲子，陳仲子處於陵者。人以爲廉，謂以不義而與之齊國必不受之。孟子以爲仲子之義，若上章所道簞食豆羹無禮則不受，萬鍾則不辯禮義而受之也。}

「人莫大焉亡親戚君臣上下。以其小者信其大者，奚可哉？」\jiazhu[1]{人當以禮義爲正。陳仲子避兄離母，不知仁義親戚上下之敘，何可以其不廉信以爲大哉？}\par \jiazhu[1]{章指言：事有輕重，行有大小，以大包小可也，以小信大，未之聞也。}

\section{桃應問瞽瞍殺人}

桃應問曰：「舜爲天子，皋陶爲士，瞽瞍殺人，則如之何？」\jiazhu[1]{桃應，孟子弟子。問皋陶爲士官，主執罪人，瞽瞍惡暴而殺人，則皋陶如何？}

孟子曰：「執之而已矣。」\jiazhu[1]{孟子曰：皋陶執之耳。}

「然則舜不禁與？」\jiazhu[1]{桃應以爲舜爲天子，使有司執其父，不禁止之邪？}

曰：「夫舜惡得而禁之？夫有所受之也。」\jiazhu[1]{夫，辭也。孟子曰：夫舜惡得禁之？夫天下乃受之於堯，當爲天理民，王法不曲，豈得禁之也？}

「然則舜如之何？」\jiazhu[1]{應問舜爲之將如何？}

曰：「舜視棄天下猶棄敝蹝也。竊負而逃，遵海濱而處，終身訢然，樂而忘天下。」\jiazhu[1]{孟子曰：舜視棄天下如捐棄敝蹝。蹝，草履可蹝者也。敝，喻不惜。舜必負父而遠逃，終身訢然，忽忘天下之爲貴也。}\par \jiazhu[1]{章指言：奉法承天，政不可枉。大孝榮父，遺棄天下，虞舜之道，趨將若此。孟子之言，揆聖意也。}

\section{居移氣，養移體}

孟子自范之齊，望見齊王之子，喟然歎曰：「居移氣，養移體，大哉居乎！夫非盡人之子與？」\jiazhu[1]{范，齊邑，王庶子所封食也。孟子之范，見王子之儀聲氣高涼，不與人同，還至齊，謂諸弟子喟然而嘆曰：「居尊則氣高，居卑則氣下，居之移人氣志，使之高涼，若供養之移人形身，使充盛也。大哉居乎者，言當慎所居，人必居仁也。凡人與王子，豈非盡是人之子也？王子居尊勢，故儀聲如是也。」}\par \jiazhu[1]{章指言：人性皆同，居使之異。君子居仁，小人處利，譬猶王子，殊於衆品也。}

孟子曰：「王子宫室、車馬、衣服多與人同，而王子若彼者，其居使之然也。況居天下之廣居者乎？」\jiazhu[1]{言王子宫室乘服皆人之所用之耳，然而王子若彼高涼者，居勢位故也。況居廣居，謂行仁義。仁義在身，不言而喻也。}

「魯君之宋，呼於垤澤之門。守者曰：『此非吾君也，何其聲之似我君也？』此無他，居相似也。」\jiazhu[1]{垤澤，宋城門名也。人君之聲相似者，以其俱居尊勢，故音氣同也。以城門不自肯夜開，故君自發聲。}\par \jiazhu[1]{章指言：輿服器用，人用不殊，尊貴居之，志氣以舒。是以居仁由義，盎然內優。胸中正者，眸子不瞀也。}

\section{食而弗愛，豕交之}

孟子曰：「食而弗愛，豕交之也；愛而不敬，獸畜之也。恭敬者，幣之未將者也。恭敬而無實，君子不可虛拘。」\jiazhu[1]{人之交接，但食之而不愛，若養豕也；愛而不敬，若人畜禽獸，但愛而不能敬也。且恭敬者，如有幣帛，當以行禮，而未以命將行之也。恭敬貴實，如其無實，何可虛拘致君子之心也？}\par \jiazhu[1]{章指言：取人之道，必以恭敬，恭敬貴實，虛則不應。實者，言敬愛也。}

\section{形色天性}

孟子曰：「形色，天性也。」\jiazhu[1]{形，謂君子體貌嚴尊也，《尚書·洪範》一曰貌。色，謂婦人妖麗之容，《詩》云「顏如蕣華」。此皆天假施於人也。}

「惟聖人然後可以踐形。」\jiazhu[1]{踐，履居之也。《易》曰：「黃中通理。」聖人內外文明，然後能以正道履居此美形。不言居色，主名尊陽抑陰之義也。}\par \jiazhu[1]{章指言：體德正容，大人所履。有表無裏，謂之柚%\KR{3031}
。是以聖人乃堪踐形也。}

\section{齊宣王欲短喪}

齊宣王欲短喪。公孫丑曰：「爲朞之喪，猶愈於已乎？」\jiazhu[1]{齊宣王以三年之喪爲太長久，欲減而短之，因公孫丑使自以其意問孟子。既不能三年喪，以朞年差愈於止而不行喪者。}

孟子曰：「是猶或紾其兄之臂，子謂之姑徐徐云爾，亦教之孝悌而已矣。」\jiazhu[1]{紾，戾也。孟子言有人戾其兄之臂，爲不順也，而子謂之曰且徐徐云爾，是豈以徐之爲差者乎？不若教之以孝悌，勿復戾其兄之臂也。今欲行其朞喪，亦猶曰徐徐之類也。}

「王子有其母死者，其傅爲之請數月之喪。公孫丑曰：若此者何如也？」\jiazhu[1]{丑曰：王之庶夫人死，迫於適夫人，不得行其喪親之數。其傅爲請之於君，欲使得行數月喪，如之何？}

曰：「是欲終之而不可得也。雖加一日愈於已，謂夫莫之禁而弗爲者也。」\jiazhu[1]{孟子曰：如是王子欲終服其子禮而不能者也。加益一日則愈於止，況數月乎？所謂不當者，謂無禁自欲短之，故譏之也。}\par \jiazhu[1]{章指言：禮斷三年，孝者欲益。富貴怠厭，思減其日。君子正言，不可阿情。丑欲朞之，故譬以紾兄徐徐也。}

\section{君子之教者五}

孟子曰：「君子之所以教者五。」\jiazhu[1]{教民之道有五品。}

「有如時雨化之者。」\jiazhu[1]{教之漸漬而沾洽也。}

「有成德者，有達財者，有答問者，有私淑艾者。」\jiazhu[1]{私，獨；淑，善；艾，治也。君子獨善其身，人法其仁，此亦與教法之道無差也。}

「此五者，君子之所以教也。」\jiazhu[1]{申言之，孟子貴重此教之道。}\par \jiazhu[1]{章指言：教人之術，莫善五者。養育英才，君子所珍。聖所不倦，其惟誨人乎？}

\section{公孫丑問道高難及}

公孫丑曰：「道則高矣，美矣，宜若登天然，似不可及也。何不使彼爲可幾及而日孳孳也？」\jiazhu[1]{丑以爲聖人之道大高遠，將若登天，人不能及也。何不少近人情，令彼凡人可庶幾，使日孳孳自勉也？}

孟子曰：「大匠不爲拙工改廢繩墨，羿不爲拙射變其彀率。君子引而不發，躍如也。中道而立，能者從之。」\jiazhu[1]{大匠不爲新學拙工故爲之改鑿廢繩墨必正也，羿不爲新學拙射者變其彀率之法也。彀，弩張嚮表率之正體，望之極思用巧之時，不可變也。君子謂於射則引弓彀弩而不發，以待彀偶也；於道則中道德之中，不以學者不能，故卑下其道，將以須於能者往取之也。}\par \jiazhu[1]{章指言：曲高和寡，道大難追。然而履正者不枉，執德者不回，故曰人能弘道。丑欲下之，非也。}

\section{天下有道以道殉身}

孟子曰：「天下有道，以道殉身；天下無道，以身殉道。未聞以道殉乎人者也。」\jiazhu[1]{殉，從也。天下有道，得行王政，道從身施功實也。天下無道，道不得行，以身從道，守道而隱。不聞以正道從俗人也。}\par \jiazhu[1]{章指言：窮達卷舒，屈伸異變，變流從顧，守者所慎，故曰金石獨止，不殉人也。}

\section{公都子問滕更}

公都子曰：「滕更之在門也，若在所禮而不答，何也？」\jiazhu[1]{滕更，滕君之弟，來學於孟子也。言國君之弟而樂在門人中，宜答見禮，而夫子不答，何也？}

孟子曰：「挾貴而問，挾賢而問，挾長而問，挾有勳勞而問，挾故而問，皆所不答也。滕更有二焉。」\jiazhu[1]{挾，接也。接己之貴勢，接己之有賢才，接己長老，接己當有功勞之恩，接己與師有故舊之好。凡恃此五者而以學問望師之待以異意而教之，皆所不當答。滕更有二焉，接貴、接賢，故不答矣。}\par \jiazhu[1]{章指言：學尚虛己，師誨貴乎。是以滕更恃二，孟子弗應。}

\section{於不可已而已者}

孟子曰：「於不可已而已者，無所不已。於所厚者薄，無所不薄也。其進銳者，其退速。」\jiazhu[1]{已，棄也。於義所不當棄而棄之，則不可，所以不可而棄之，使無罪者咸恐懼也。於義當厚而反薄之，何不薄也？不憂見薄者亦皆自安矣。不審察人而過進不肖，越其倫，悔而退之必速矣。當翔而後集，慎如之何？}\par \jiazhu[1]{章指言：賞僭及淫，刑濫傷善，不僭不濫，詩人所紀。是以季文三思，何後之有？}

\section{君子之於物}

孟子曰：「君子之於物也，愛之而弗仁。」\jiazhu[1]{物，謂凡物可以養人者也，當愛育之而不知人仁，若犧牲不得不殺也。}

「於民也，仁之而弗親。」\jiazhu[1]{臨民以非己族類，故不得與親同也。}

「親親而仁民，仁民而愛物。」\jiazhu[1]{先視其親戚，然後仁民，仁民然後愛物，用恩之次也。}\par \jiazhu[1]{章指言：君子布德，各有所施，事得其宜，故謂之義也。}

\section{知者無不知，當務之急}

孟子曰：「知者無不知也，當務之爲急；仁者無不愛也，急親賢之爲務。」\jiazhu[1]{知者知所務善也，仁者務愛賢也。}

「堯舜之知而不遍物，急先務也；堯舜之仁不遍愛人，急親賢也。」\jiazhu[1]{物，事也。堯舜不遍知百工之事，不遍愛衆人，先愛賢，使治民，不二三自往親加恩惠也。}

「不能三年之喪，而緦小功之察；放飯流歠，而問無齒決，是之謂不知務。」\jiazhu[1]{尚不能行三年之喪，而復察緦麻小功之禮。放飯，大飯也；流歠，長歠也。齒決，斷肉置其餘也。於尊者前賜食，大飯長歠，不敬之大者；齒決，小過耳。言世之先務，舍大譏小，若此之類也。}\par \jiazhu[1]{章指言：振裘持領，正羅維綱。君子百行，先務其崇。是以堯舜親賢，大化以隆，道爲要也。}
